% PrototypingAndWireframes.tex

\raggedright

\section{Overview of Prototyping}

The prototyping stage was used to move TAC from an initial project idea into a working system design. This stage was important because the project involved both software development and infrastructure implementation. The system could not be designed only on paper, because several parts of the project depended on practical testing, including web application layout, database structure, Active Directory integration, ADFS authentication, firewall rules, VLAN separation, SSL certificate behaviour, and SQL Server connectivity.

The purpose of prototyping was to test ideas early before committing to the final implementation. For the software side, this involved creating early versions of the user interface, testing how case and subject records would be displayed, and using mock data before the SQL database was fully connected. For the infrastructure side, prototyping involved building and testing individual services such as pfSense, Active Directory, ADFS, SQL Server, and the Linux web server before joining them into the final architecture.

The project followed an iterative approach. Early versions were intentionally simple and focused on proving one part of the design at a time. For example, the first version of the web application used local mock data so that the layout, navigation, and user workflows could be tested quickly. Once the interface and record structure became clearer, the project moved toward SQL backed API routes. This allowed the system to evolve from a visual prototype into a persistent application.

Prototyping also helped identify technical problems early. Issues such as DNS resolution, domain trust behaviour, SQL connectivity, ADFS configuration, and certificate trust needed practical testing. These problems would have been difficult to fully predict from the design alone. By building the system in stages, each problem could be isolated and resolved before the next layer was added.

% Suggested figure:
% Figure 5.1: Prototype evolution diagram showing idea, mock UI, SQL integration, identity integration, and final platform.
% Caption: TAC prototype development lifecycle.

\section{Software Prototyping Approach}

The software prototype began with the main goal of creating a usable case management interface. Before connecting the system to SQL Server or Active Directory, the first priority was to define how users would move through the application and how the main records would be presented.

The early application prototype focused on the dashboard layout, navigation sidebar, case registry, subject registry, analytics page, and audit page. These areas formed the core structure of the final web application. A dashboard style layout was selected because it is familiar in administrative and operational systems. It allows users to move between different functions without leaving the main application shell.

The first prototype used mock data. This allowed the interface to be tested without waiting for the database schema and API routes to be completed. Mock data was useful because it allowed realistic case and subject records to be displayed in tables, cards, and detail views. This helped test whether the information was readable, whether the layout was clear, and whether the user interface could handle multiple records.

Once the mock interface worked, the prototype was refined. Case records were given clearer fields, including case type, status, priority, sensitivity level, department, dates, and descriptive notes. Subject records were also refined so that they could represent different categories such as suspect, victim, witness, reporter, missing person, or other relevant subject types.

The software prototype also helped shape the final terminology. The system originally used the term person, but this was later changed to subject. This was a more appropriate term for an investigative style system because it is broader and avoids implying that every record has the same role or legal meaning.

The final software prototype provided the foundation for the implementation phase. Once the interface and workflows were stable, development moved toward API integration, SQL persistence, audit logging, and role aware behaviour.

% Suggested figure:
% Figure 5.2: Screenshot of early case registry prototype using mock data.
% Caption: Early software prototype showing case registry layout.

\section{Infrastructure Prototyping Approach}

The infrastructure prototype was developed alongside the software prototype. This was necessary because TAC was not designed as a standalone application. The final platform depended on several infrastructure services working together correctly.

The infrastructure prototype began with the virtual environment. Proxmox was used to host the required virtual machines and separate the different system roles. This allowed systems such as Windows Server, SQL Server, Linux web hosting, pfSense, and monitoring services to be built and tested independently.

The network prototype focused on VLAN separation and pfSense firewalling. Before the application could be treated as secure, the network paths between users, servers, and management systems needed to be understood. pfSense was used to test routing between VLANs and to define which services should be reachable between network zones.

The identity prototype focused on Active Directory, DNS, domain trust, and ADFS. Active Directory was tested first because it provided the identity foundation. DNS was then validated because domain authentication and ADFS depend heavily on correct name resolution. ADFS was tested later as the federation layer for the application.

The database prototype focused on SQL Server connectivity from the Linux web application server. This was an important milestone because the final application needed to move from mock data to persistent SQL-backed data. Testing confirmed that the web server could connect to SQL Server over the internal network when firewall rules, SQL configuration, and connection details were correct.

The infrastructure prototype also included SSL/TLS testing. Internal HTTPS access was required so that the platform could demonstrate encrypted communication. Certificate configuration was therefore tested for the web application and identity services.

This infrastructure prototyping approach helped reduce risk. Instead of attempting to deploy every component at once, each layer was built, tested, and then connected to the wider architecture.

% Suggested figure:
% Figure 5.3: Infrastructure prototype diagram showing staged setup of pfSense, AD, ADFS, SQL Server, Linux web app, and Wazuh.
% Caption: Infrastructure prototyping stages used during TAC development.

\section{Initial Mock Data Prototype}

The initial mock data prototype was one of the most useful stages in the software development process. At the start of the project, the database schema and SQL Server connection were not yet finalised. However, the application still needed to be designed and tested from a user perspective. Mock data allowed this to happen quickly.

Mock case records were used to populate the case registry. These records included realistic fields such as case title, type, status, priority, sensitivity, department, and date information. This allowed the table layout, filters, cards, and detail views to be tested before live database records were available.

Mock subject records were used in a similar way. These records helped test how subjects would be displayed and how the user interface would handle different subject categories. Using mock subjects also helped identify the need for consistent formatting. For example, enum style values with underscores were later improved so that they displayed more clearly to users.

The mock data prototype also supported early analytics. Even before SQL Server integration, sample records allowed the dashboard to show case counts, status groupings, and other summary information. This helped prove that the analytics area had value and should remain part of the final system.

A major advantage of the mock data stage was speed. Interface changes could be made quickly without needing to change the database each time. This allowed the layout to mature before more complex backend work was added.

However, mock data was only suitable for early prototyping. It did not provide real persistence, multi user behaviour, database validation, or realistic backend integration. For this reason, the project later moved from mock data to SQL backed API routes. This transition was important because it changed the system from a visual prototype into a working application with persistent storage.

% Suggested figure:
% Figure 5.4: Mock data example displayed in the case registry.
% Caption: Mock data prototype used to validate case registry layout.

\section{User Interface Wireframes}


Wireframes were used to plan the structure of the TAC user interface before and during implementation. The aim of the wireframes was not to create a perfect visual design at the beginning, but to decide how users would move through the system and where the main features would be placed.

The main design decision was to use a dashboard shell. This shell included a sidebar navigation area and a main content area. The sidebar provided access to the key application sections: cases, subjects, analytics, and audit logs. The main content area changed depending on the selected page.

The wireframes focused on clarity and consistency. Each major page followed a similar pattern: a page heading, a short description, action buttons where appropriate, and the main content displayed in tables, cards, forms, or dashboards. This consistency was important because the platform needed to feel like one connected system rather than a collection of unrelated pages.

The case registry wireframe was designed around a list of case records. It needed to show enough information for users to identify a case quickly without overwhelming the screen. Important fields such as case reference, title, status, priority, type, and sensitivity were prioritised.

The subject registry wireframe followed a similar structure. It needed to show subject identity information, subject type, confidence level, and related details. The design also needed to allow users to open individual subject records for more detailed information.

The analytics wireframe focused on summary information. It was designed to show high level indicators such as total cases, total subjects, case status breakdowns, and priority distribution. This helped demonstrate that the platform could support oversight as well as record management.

The audit wireframe focused on accountability. It needed to show recorded actions in a way that could be reviewed easily, including action type, timestamp, and affected record.

% Suggested figure:
% Figure 5.5: Wireframe of the TAC dashboard shell.
% Caption: Dashboard shell wireframe used to structure the TAC interface.

\section{Case Registry Wireframe}

The case registry wireframe was one of the most important interface designs because case management is the central feature of TAC. The wireframe needed to support fast recognition of records while still providing enough detail for operational use.

The case registry was designed as a structured list view. Each case entry needed to display a unique case reference, title, case type, status, priority, sensitivity level, and department. These fields were selected because they allow users to quickly understand the nature and importance of a case.

Action controls were also considered in the wireframe. Users with the correct permissions would need options to create a new case, open case details, edit a case, or delete a record. However, these actions should not be available equally to every role. This meant that the interface needed to support role-aware behaviour later in the implementation.

The wireframe also considered visual hierarchy. Critical or high priority information should be easy to identify. Case status and priority were therefore suitable for badge style visual treatment. This makes the registry easier to scan than plain text alone.

The case detail view was designed as a more focused page. Instead of showing many cases, it shows one selected case with full information. This page allows the user to review the case description, metadata, dates, and related fields in more detail.

The case registry wireframe evolved during the project. Early versions displayed some information inside and outside the page shell, which was later refined so that content appeared consistently inside the application layout. This improved the professional appearance and usability of the final interface.

% Suggested figure:
% Figure 5.6: Case registry wireframe or final screenshot.
% Caption: Case registry layout used for managing TAC case records.

\section{Subject Registry Wireframe}

The subject registry wireframe was designed to manage records relating to individuals or entities connected to case work. The term subject was chosen because it is more suitable for an investigative case management environment than the more general term person.

The subject registry needed to display key identifying and classification information. This included subject name, subject type, identity confidence, contact or reference information where appropriate, and any relevant classification fields. The wireframe needed to allow users to identify subjects quickly while avoiding unnecessary clutter.

The subject registry followed a similar layout pattern to the case registry. This consistency was intentional. Users should not need to learn a completely different interface when moving between cases and subjects. The same general structure of heading, actions, list view, and detail access was used.

The subject detail page was designed to show more complete information about a selected subject. This could include personal or descriptive details, classification information, notes, and possible links to case records. In a real environment, this area would require strong access control because subject records may contain sensitive personal data.

The wireframe also supported future linking between cases and subjects. A subject may be connected to more than one case, and a case may involve more than one subject. Even where advanced relationship mapping was not the main focus of the final system, the interface was designed with this future expansion in mind.

Like the case registry, the subject registry was refined during implementation. The final design aimed to keep the interface clean, consistent, and aligned with the rest of the platform.

% Suggested figure:
% Figure 5.7: Subject registry wireframe or final screenshot.
% Caption: Subject registry layout used for managing subject records.

\section{Bulk Import Wireframe}

The bulk import wireframe was designed to support the ingestion of multiple synthetic records into the system. This feature was included because manually entering every test case and subject would be inefficient and would limit the scale of testing. Bulk import allowed the platform to be populated with enough data to demonstrate listing, analytics, audit logging, and database persistence.

The wireframe for bulk import needed to be simple and controlled. Users should be able to select an import file, submit it, and receive feedback on whether the import succeeded or failed. The page also needed to explain what type of data structure was expected so that records could be mapped correctly into the system.

Bulk import was especially useful for project validation. It allowed synthetic case data to be generated and imported for testing. This made the final demonstration stronger because the system could show realistic volumes of records rather than only a small number of manually created examples.

The bulk import design also needed to consider errors. If a file was invalid, missing required fields, or contained unsuitable data, the system should not fail silently. The user should receive useful feedback. This requirement influenced the implementation of validation and error handling around the import process.

From an ethical perspective, the bulk import feature was tested using synthetic data only. No real personal data, real case records, victim data, suspect data, or operational law enforcement data was used. This allowed the feature to be demonstrated safely while avoiding data protection concerns.

The wireframe therefore supported both functionality and responsible testing. It helped show that the system could ingest structured records while keeping the testing process controlled and ethical.

% Suggested figure:
% Figure 5.8: Bulk import page wireframe or screenshot.
% Caption: Bulk import interface used to upload synthetic case and subject data.

\section{Analytics Dashboard Wireframe}

The analytics dashboard wireframe was designed to give users a high-level view of the data stored in the platform. While the case and subject registries focus on individual records, the analytics dashboard focuses on summary information.

The dashboard was designed to display key indicators such as total number of cases, total number of subjects, open cases, closed cases, priority distribution, and case status breakdowns. These summaries help users understand the overall state of the dataset without manually reviewing every record.

The wireframe used cards and chart areas to separate different types of information. Summary cards are useful for quick totals, while charts are useful for comparing categories such as status or priority. This combination improves readability and makes the dashboard more useful during demonstrations.

The analytics dashboard was also important for testing the data model. Analytics only become meaningful when records are stored in a structured and consistent way. Therefore, the dashboard helped validate that case and subject fields were suitable for grouping and counting.

The dashboard also benefited from the bulk import feature. A small number of manual records would not show meaningful patterns. By importing synthetic data, the analytics dashboard could demonstrate more realistic behaviour.

From a design perspective, the analytics page needed to remain simple. The aim was not to create an advanced intelligence analysis platform, but to demonstrate that the case management system could provide useful oversight based on stored records.

% Suggested figure:
% Figure 5.9: Analytics dashboard wireframe or screenshot.
% Caption: Analytics dashboard used to summarise TAC case and subject records.

\section{Audit Log Wireframe}

The audit log wireframe was designed to support accountability. In a case management system, users should not be able to create, change, import, or delete records without leaving a trace. The audit log provides a way to review important activity.

The audit page was designed as a structured list of events. Each event should show information such as the action performed, the affected record, the timestamp, and the user or system identity where available. This format makes audit activity easier to review and supports later validation.

The audit log wireframe needed to be readable and filterable in future versions. While the final implementation focused on core audit visibility, the layout was designed so that future filtering by user, action type, date, or record could be added.

The audit log also supports security and data protection goals. If sensitive records are handled, the organisation should be able to demonstrate accountability. Audit logs help provide this accountability by showing what happened inside the system.

The audit wireframe also connects to the wider monitoring design. Application-level audit logs show actions inside TAC, while Wazuh provides infrastructure-level monitoring. These two layers complement each other. One explains what users did in the application, while the other helps observe what is happening on the underlying systems.

The audit log page was therefore included as a core part of the system rather than an optional administrative extra. It supports the overall project aim of building a secure and accountable case management platform.

% Suggested figure:
% Figure 5.10: Audit log wireframe or final screenshot.
% Caption: Audit log layout used to review important TAC user actions.

\section{Prototype Feedback and Iteration}

The prototype changed several times during development. These changes were based on testing, usability review, technical limitations, and the need to keep the final scope achievable.

One major change was the move from mock data to SQL backed storage. The mock data prototype was useful for early interface design, but it was not suitable for the final system. SQL Server integration became necessary so that records could persist and the system could behave more like a real enterprise platform.

Another change was the refinement of terminology. The interface moved from person-based language toward subject-based language. This made the system sound more suitable for investigative case management and reduced ambiguity.

The layout was also improved during prototyping. Some early pages displayed content inconsistently inside and outside the application shell. This was later corrected so that pages followed the same visual structure. This made the application feel more complete and professional.

The scope was also refined during development. Some ideas that were considered during the design phase were not included in the final implementation due to time constraints. These decisions were part of the agile development process. The project prioritised the core features that were most important to the final system: case management, subject management, SQL persistence, identity integration, RBAC design, VLAN segmentation, audit logging, and monitoring.

This iterative process improved the final platform. Rather than attempting to build every possible feature, the project focused on delivering a coherent and demonstrable system. The prototype feedback cycle helped identify what should be kept, refined, deferred, or removed.

% Suggested figure:
% Figure 5.11: Iteration timeline showing mock prototype, UI refinement, SQL migration, RBAC, and validation.
% Caption: Iterative prototype refinement during TAC development.


\section{Use of AI Tools During Prototyping}

AI tools were used as supporting tools during the prototyping stage. Their use was limited to assistance, exploration, and acceleration rather than replacing the project work. The final design decisions, implementation choices, infrastructure configuration, testing, and report structure were reviewed and adapted by the author.

Claude was referenced during the user interface prototyping stage. It was used as a source of inspiration for layout ideas, interface wording, and component organisation. The final interface was not copied directly from a generated design. Instead, AI-assisted suggestions were adapted to fit the requirements, project scope, and existing codebase.

ChatGPT was used to support the generation of synthetic test data and to help plan bulk import scenarios. This was useful because the platform needed realistic looking records for testing, but it would have been inappropriate to use real personal or investigative data. Synthetic data made it possible to test import workflows, analytics, audit logging, and database behaviour while avoiding ethical and data protection concerns.

AI support was also used during problem solving and documentation planning. For example, it helped structure explanations, identify testing categories, and organise project sections. However, the technical implementation remained dependent on practical configuration, code changes, testing, screenshots, and validation carried out in the project environment.

The use of AI tools was therefore transparent and controlled. AI was used as an assistant during prototyping and documentation support, not as a substitute for understanding or implementation. This aligns with the requirement to acknowledge AI usage openly and ensures that the work remains accountable.

% Suggested table:
% Table 5.1: AI tool usage during prototyping.
% Columns: Tool, Purpose, Type of Support, Final Author Control.

\section{Prototyping Summary}

The prototyping stage was essential to the successful development of TAC. It allowed the project to move from an initial concept into a clearer working design before final implementation. Both software and infrastructure prototypes were required because the final system depended on the interaction between the web application, database, identity services, network segmentation, certificates, and monitoring.

The software prototype helped define the user interface, dashboard layout, case registry, subject registry, bulk import page, analytics dashboard, and audit log. Mock data allowed the early interface to be developed quickly and then replaced later with SQL backed persistent storage.

The infrastructure prototype helped validate the technical environment. Proxmox, pfSense, VLANs, Active Directory, ADFS, SQL Server, Linux hosting, SSL/TLS, and Wazuh were tested in stages. This reduced risk and allowed problems to be solved before the final system was assembled.

The prototyping process also supported good project management. Features were refined, renamed, deferred, or prioritised based on their value and feasibility. This helped ensure that the final implementation focused on the most important parts of the platform.

Overall, prototyping made the final system stronger. It improved usability, reduced technical risk, supported ethical testing through synthetic data, and provided a clear path from concept to implementation.

% Suggested figure:
% Figure 5.12: Final prototype-to-implementation summary diagram.
% Caption: Summary of how the TAC prototype evolved into the final implementation.